% Figure — Oscillateur à pont réactif (a) et montage réactif seul (b) du TD 6 (EANA3, chapitre 4).
% D'après le support [support, p. 24, figures 1 et 2], redessiné.
% Rendu SVG :
%   latex -output-directory=<dossier temporaire> 04-pont-reactif.tex
%   dvisvgm --no-fonts --exact-bbox -o 04-pont-reactif.svg <dossier temporaire>/04-pont-reactif.dvi
\documentclass[border=4pt]{standalone}
\usepackage[T1]{fontenc}
\usepackage[utf8]{inputenc}
\usepackage{amsmath}
\usepackage{sansmath}
\renewcommand{\familydefault}{\sfdefault}
\usepackage[european, straightvoltages, americancurrents]{circuitikz}
\usetikzlibrary{shapes.geometric,calc,positioning}
\begin{document}
\sansmath
\begin{circuitikz}[font=\small, >=latex, fleche/.style={->, thick}]
  % (b) montage réactif
  \begin{scope}[xshift=7.4cm]
    \node at (-0.4,3.6) {(b)};
    \draw (2.0,1.6) node[op amp, noinv input up] (oa) {};
    \draw (0,2.8) node[ocirc]{} -- (0.9,2.8) coordinate (x) -- (x |- oa.+) -- (oa.+);
    \draw (x) to[R, l=$R_f$] (3.6,2.8) -- (3.6,0 |- oa.out) -- (oa.out);
    \draw (3.6,0 |- oa.out) to[R, l=$R$] (3.6,0.3) coordinate (m) to[R, l=$R$] (3.6,-1.2) node[ground]{};
    \draw (m) -- (0.9,0.3) -- (0.9,0 |- oa.-) -- (oa.-);
    \draw (0,-1.2) node[ocirc]{};
    \draw[fleche] (0.2,3.05) -- node[above]{$I_{in}$} (0.7,3.05);
    \draw[fleche] (-0.3,-1.0) -- node[left]{$V_{in}$} (-0.3,2.6);
    \draw[fleche] (1.4,3.35) -- node[above]{$V_f$} (3.2,3.35);
  \end{scope}
  % (a) oscillateur complet
  \node at (-0.4,3.6) {(a)};
  \draw (3.6,1.6) node[op amp, noinv input up] (ob) {};
  \draw (1.8,2.8) coordinate (y) -- (y |- ob.+) -- (ob.+);
  \draw (y) to[R, l=$R_f$] (5.2,2.8) -- (5.2,0 |- ob.out) -- (ob.out);
  \draw (5.2,0 |- ob.out) to[R, l=$R$] (5.2,0.3) coordinate (m2) to[R, l=$R$] (5.2,-1.2) node[ground]{};
  \draw (m2) -- (2.5,0.3) -- (2.5,0 |- ob.-) -- (ob.-);
  \draw (y) -- (1.8,1.2) coordinate (h);
  \draw (h) -- (-0.2,1.2);
  \draw (-0.2,1.2) to[cute inductor, l_=$L$] (-0.2,-1.2);
  \draw (0.6,1.2) to[C, l=$C$] (0.6,-1.2);
  \draw (1.4,1.2) to[R, l=$R_p$] (1.4,-1.2);
  \draw (-0.2,-1.2) -- (1.8,-1.2) node[ground]{};
  \draw (1.8,1.2) -- (1.4,1.2);
\end{circuitikz}
\end{document}
